\section{Preliminaries}
\label{sec:prelim}

We fix notation and recall standard definitions from universal algebra and term
rewriting theory.  The reader may consult Baader and Nipkow \cite{baader1999term}
for a comprehensive treatment.

\subsection{Magmas and Equational Laws}

\begin{definition}
A \emph{magma} is a pair $(M, *)$ where $M$ is a nonempty set and
$* \colon M \times M \to M$ is a binary operation.  We impose no additional axioms.
\end{definition}

Let $\mathcal{X} = \{x, y, z, w, \ldots\}$ be a countably infinite set of
\emph{variables} and let $\Sigma = \{*\}$ be the signature consisting of a single
binary operation symbol.  The set $T(\Sigma, \mathcal{X})$ of \emph{terms} over
$\Sigma$ and $\mathcal{X}$ is defined inductively: every variable $x \in \mathcal{X}$
is a term, and if $s, t$ are terms then $s * t$ is a term.

\begin{definition}
An \emph{equational law} is a pair of terms $(s, t) \in T(\Sigma, \mathcal{X})^2$,
written $s = t$.  An equational law $s = t$ is \emph{trivial} if $s$ and $t$ are
identical as terms.
\end{definition}

\begin{definition}
A magma $(M, *)$ \emph{satisfies} an equational law $s = t$, written
$M \models (s = t)$, if for every assignment $\varphi \colon \mathcal{X} \to M$ of
variables to elements, the evaluations $\varphi(s)$ and $\varphi(t)$ in $M$ coincide.
Here $\varphi$ extends to all terms by $\varphi(u * v) = \varphi(u) * \varphi(v)$.
\end{definition}

\begin{definition}\label{def:implication}
An equational law $E$ \emph{implies} an equational law $E'$, written $E \models E'$,
if every magma satisfying $E$ also satisfies $E'$.  We write $E \not\models E'$ for
the negation; in this case any magma satisfying $E$ but not $E'$ is a
\emph{counterexample} to the implication.
\end{definition}

\begin{example}
Let $E$ be the \emph{left-zero} law $x * y = x$ and let $E'$ be $x * y = y * x$
(commutativity).  Then $E \models E'$: if $a * b = a$ for all $a, b$ then
$a * b = a = b * a$ for all $a, b$.  Conversely, let $E$ be commutativity.  The
magma on $\{0, 1\}$ with $0*0=1, 0*1=1, 1*0=1, 1*1=0$ satisfies commutativity
but not left-zero, so commutativity does not imply left-zero.
\end{example}

\begin{remark}
The implication relation $\models$ is a preorder on equational laws.  The trivial
law $x = x$ is satisfied by every magma and hence is implied by every law.  The
law $x = y$ (identifying all elements) is satisfied only by singleton magmas and
implies every other law.
\end{remark}

\subsection{Term Rewriting Systems}

\begin{definition}
A \emph{rewrite rule} is an ordered pair $(\ell, r)$ of terms written $\ell \to r$.
A \emph{term rewriting system} (TRS) is a finite set $R$ of rewrite rules.
\end{definition}

For a rule $\ell \to r$, let $\operatorname{vars}(\ell)$ and $\operatorname{vars}(r)$ denote the sets of
variables occurring in $\ell$ and $r$ respectively.  A rule $\ell \to r$ is
\emph{left-linear} if each variable of $\operatorname{vars}(\ell)$ occurs at most once in $\ell$.

\begin{definition}\label{def:valid-rule}
A rewrite rule $\ell \to r$ is \emph{valid} (or \emph{regular}) if
$\operatorname{vars}(r) \subseteq \operatorname{vars}(\ell)$.  A TRS is valid if every rule is valid.
\end{definition}

The validity condition ensures that no RHS variable is ``free'' (unbound by the LHS
match), a necessary condition for the rule to be applicable under a reduction order.

\begin{definition}
Given a TRS $R$ and terms $s, t$, we write $s \to_R t$ if there exists a rule
$\ell \to r \in R$, a position $p$ in $s$, and a substitution $\sigma$ such that
$s|_p = \ell\sigma$ and $t = s[r\sigma]_p$ (replacing the subterm at position $p$
by $r\sigma$).  We write $\to_R^*$ for the reflexive-transitive closure.
\end{definition}

\begin{definition}
A term $t$ is a \emph{normal form} under $R$ if no rule of $R$ applies to any
subterm of $t$.  If every term has a normal form, $R$ is called \emph{terminating}
(or \emph{strongly normalizing}).  We write $\mathrm{NF}_R(t)$ for the normal
form of $t$ when it is unique.
\end{definition}

\begin{definition}
$R$ is \emph{confluent} if whenever $s \to_R^* t_1$ and $s \to_R^* t_2$ there
exists $u$ with $t_1 \to_R^* u$ and $t_2 \to_R^* u$.  $R$ is \emph{convergent}
if it is both terminating and confluent; in this case every term has a unique normal
form.
\end{definition}

\begin{definition}
A \emph{reduction order} $>$ on terms is a strict order that is well-founded,
stable under substitution ($s > t \Rightarrow s\sigma > t\sigma$ for all $\sigma$),
and compatible with context ($s > t \Rightarrow u[s]_p > u[t]_p$ for all contexts).
\end{definition}

\begin{definition}[Lexicographic Path Order \cite{dershowitz1987orderings}]
Fix a strict total order (precedence) on function symbols.  The \emph{Lexicographic
Path Order} (LPO) $>_{\mathrm{lpo}}$ is defined recursively: $f(s_1,\ldots,s_m)
>_{\mathrm{lpo}} g(t_1, \ldots, t_n)$ if (i) $f > g$ and $f(s_1,\ldots,s_m) >_{\mathrm{lpo}} t_j$
for all $j$, or (ii) $f = g$ and $(s_1,\ldots,s_m) >_{\mathrm{lex}} (t_1,\ldots,t_n)$
and $s_i >_{\mathrm{lpo}} t_j$ for all $j$ beyond the first differing position,
or (iii) $s_i \geq_{\mathrm{lpo}} g(t_1,\ldots,t_n)$ for some $i$.  LPO is a
reduction order.
\end{definition}

\subsection{Knuth-Bendix Completion}

\begin{definition}
Given rules $\ell_1 \to r_1$ and $\ell_2 \to r_2$ (renamed to have disjoint
variables), a \emph{critical pair} arises when a non-variable subterm $\ell_1|_p$
unifies with $\ell_2$ via most general unifier $\sigma$.  The critical pair is
$(r_1\sigma,\; \ell_1\sigma[r_2\sigma]_p)$: the two ways of rewriting $\ell_1\sigma$.
\end{definition}

The following classical result underpins the soundness of KB completion.

\begin{theorem}[Critical Pair Lemma, Huet \cite{huet1980confluent}]\label{thm:cpl}
A terminating TRS $R$ is confluent if and only if every critical pair $(s, t)$ of
$R$ is \emph{joinable}: there exists $u$ with $s \to_R^* u$ and $t \to_R^* u$.
\end{theorem}

\begin{definition}[Knuth-Bendix Completion \cite{knuth1970simple}]
\label{def:kb}
Given an equational system $E = \{E_1, \ldots, E_k\}$ and a reduction order $>$,
\emph{Knuth-Bendix completion} is a procedure that attempts to construct a convergent
TRS $R$ equationally equivalent to $E$.  At each step, equations in $E$ are oriented
into rules using $>$, critical pairs of $R$ are computed, and unjoinable pairs are
added back to $E$ as new equations.  The procedure
\begin{itemize}[leftmargin=*,itemsep=1pt,topsep=2pt]
  \item \emph{succeeds} if $E$ becomes empty (all critical pairs are joinable),
        yielding a convergent $R$;
  \item \emph{fails} if some equation cannot be oriented by $>$; or
  \item \emph{diverges} if the process never terminates.
\end{itemize}
\end{definition}

The central soundness result for KB completion is the following.

\begin{theorem}[KB Soundness \cite{knuth1970simple,baader1999term}]\label{thm:kb-classical}
Let $E$ be an equational system.  If Knuth-Bendix completion terminates successfully
with convergent TRS $R$, then for all terms $s, t$:
\[
  E \models (s = t)
  \quad\Longleftrightarrow\quad
  \mathrm{NF}_R(s) = \mathrm{NF}_R(t).
\]
\end{theorem}

\begin{proof}
By Birkhoff's completeness theorem \cite{birkhoff1935structure}, $E \models (s = t)$
iff $s \leftrightarrow_E^* t$ (the terms are connected by a sequence of equational
rewrites from $E$).  Since $R$ is equationally equivalent to $E$, we have
$s \leftrightarrow_E^* t$ iff $s \leftrightarrow_R^* t$.  Since $R$ is convergent,
$s \leftrightarrow_R^* t$ iff $\mathrm{NF}_R(s) = \mathrm{NF}_R(t)$ (both sides
reduce to the unique normal form of their equivalence class).
\end{proof}

\subsection{Finite Model Finding and SMT}

\begin{definition}
A \emph{finite magma counterexample} of size $n$ for the implication $E \models E'$
is a magma $(M, *)$ with $|M| = n$ such that $M \models E$ but $M \not\models E'$.
\end{definition}

\begin{proposition}[Finite Model Finding Completeness]\label{prop:z3-complete}
For any fixed $n \geq 1$, there is a decision procedure that determines in finite
time whether a size-$n$ magma counterexample to $E \models E'$ exists.
\end{proposition}

\begin{proof}
A magma of size $n$ is determined by its multiplication table: an $n \times n$
array $T$ with $T_{ij} \in \{0,\ldots,n-1\}$.  The constraint ``$M \models E$''
for $E = (\ell_E = r_E)$ is a conjunction of $n^{|\operatorname{vars}(E)|}$ equalities, one
per variable assignment, each expressible as an arithmetic constraint on the
$T_{ij}$.  Similarly, ``$M \not\models E'$'' is a disjunction of $n^{|\operatorname{vars}(E')|}$
disequalities.  The resulting propositional arithmetic formula is finite and can be
decided by any complete SAT/SMT solver (for instance, by exhaustive search or by
the DPLL(T) algorithm \cite{demoura2008z3}).
\end{proof}

\begin{remark}
Proposition~\ref{prop:z3-complete} establishes completeness for each \emph{fixed}
$n$.  The implication problem over \emph{all} finite magmas---asking whether $E$
and $E'$ disagree on some finite magma---remains undecidable in general
\cite{mckenzie1971spectra}.  Moreover, by Austin's theorem \cite{austin1979equational},
finite implication ($E \models_\mathrm{fin} E'$: every finite magma satisfying $E$
also satisfies $E'$) is strictly weaker than the full implication $E \models E'$.
\end{remark}
